\section{Introduction}
\label{sec:intro}

An analytics agent pointed at a warehouse it does not understand fails in a
characteristic way. It reads the schema, writes plausible SQL, and returns a
number that is wrong for a reason no column name records: a code whose values
mean something specific to the business, a join whose cardinality depends on a
convention, a \texttt{NULL} that means \emph{applies to everything} rather than
\emph{unknown}. The query runs. The answer is confidently wrong.

The industry's answer is the \emph{context layer}: curated documentation
stored alongside the data and fetched by the agent at query time. The reported
gains are large. MotherDuck's Guides raise DABStep accuracy by 72 percentage
points while cutting cost by roughly 55\%~\cite{motherduck-guides}, and
similar claims accompany every semantic-layer and context-layer product now
shipping.

These results share a shape: system with the layer against system without
it, two accuracy numbers, a large difference. That establishes \emph{that}
context layers work. It cannot establish \emph{which part} works, and the part
matters because the parts cost different amounts. A context layer as deployed
bundles at least three things:

\begin{enumerate}
\itemsep2pt
\item \textbf{Semantic content}: prose describing what the columns mean, how
      the metrics are defined and which conventions hold.
\item \textbf{Retrieval scaffolding}: the tools that fetch that content, the
      instruction to fetch it before writing SQL, and the narrowed table
      surface that comes with a curated catalogue.
\item \textbf{Pre-computed artifacts}: views and macros that encode a metric
      once, so the agent calls it instead of deriving it.
\end{enumerate}

Semantic content is written once per domain and often exists already, in a
vendor manual or a data dictionary. Scaffolding is engineering, paid once.
Pre-computed macros are paid \emph{per metric}, forever, and cover only the
metrics somebody anticipated. A practitioner deciding where to spend needs to
know which of the three bought the 72 points.

Prior work has isolated the first. BIRD withholds and supplies human-written
evidence, worth 20 points to GPT-4~\cite{bird}; KaggleDBQA reports that
database documentation doubles parser accuracy~\cite{kaggledbqa}; an ontology
layer over an enterprise schema moves GPT-4 from 16\% to
54\%~\cite{sequeda-kg}. None separates content from the scaffolding that
delivers it, and none says what pre-computing the same content would have
bought.

\subsection*{Contribution}

We report a controlled ablation on DABStep~\cite{dabstep}, run on four model
families, that separates all three. The instrument is a \emph{data contract}:
a YAML artifact that carries a domain's semantics (what a fee is, how a rule
matches a transaction) together with the rules an agent's tools enforce
(which tables may be read, which operations are forbidden). Two properties of
that artifact make the separation possible.

First, the contract can be \emph{hollowed}. A generator empties every field
of prose --- domain summaries, metric descriptions, the SQL expression that
carries the fee formula, column and relationship descriptions --- and leaves
names, types, structure, the table allow-list, the forbidden operations and
all nine governed tools byte-for-byte identical. An agent given the hollow
contract still receives the instruction to look things up, still has the
tool, and still calls it. It is simply told nothing when it does.

Second, the contract \emph{compiles}. Its declared SQL expressions, transcribed
into two views, reproduce the benchmark's gold answer on every one of the 176
tasks they cover. That is the accuracy a pre-computed layer built from the
same content would reach, and it turns each arm's accuracy into a recovery
rate against a proven ceiling.

Together the two place the arms on a ladder: schema alone, scaffolding
without content, content delivered through scaffolding, and content
compiled. Beside the ladder sits the practitioner's baseline, the same
knowledge pasted verbatim into the prompt: content without scaffolding.
Reading the ladder gives four findings and one non-finding.

\textbf{Content dominates; scaffolding is real but secondary}
(Section~\ref{sec:decomposition}). Restoring the prose to the hollow contract
raises hard-task accuracy by 26 to 36 points on every model. Adding the empty
scaffolding to a bare schema is worth 0 to 5 points on the two flash models
(not significant over all tasks) and 14 to 15 on the two frontier models. The content step
exceeds the scaffolding step by $1.8{\times}$ to $6.6{\times}$, and on one
model is the whole effect.

\textbf{The contract beats the prompt because its rule is reachable, not
because it knows more} (Section~\ref{sec:why}). Both the contract arm and the
prompt arm are given the fee-matching rule; only the contract arm applies it.
Read in the submitted SQL, the contract arm writes every one of the six fee
clauses a task needs 39\%, 65\%, 55\% and 98\% of the time; the prompt arm,
3\%, 9\%, 4\% and 4\%. On one model the 24{,}000-character prompt is
indistinguishable from a bare schema on hard tasks, while the contract
carrying the same knowledge sits 45 points above both. The contract arm also
does the least work on every model, and is the cheapest way to buy a correct
answer on three of the four.

\textbf{The derivation gap closes with capability; the composition gap does
not} (Section~\ref{sec:ceiling}). On the 176 tasks the compiled contract
answers outright, the contract arm falls short of the ceiling by 39.2, 37.3,
26.1 and 5.1 points in order of model capability. On tasks that compose the
same semantics further, the shortfall moves between 54 and 40 points with no
clean ordering. The cost of making an agent derive a query rather than call a
macro is transient; the cost of the reasoning built on top of it is not yet.

\textbf{Declared rules deter mutation without being enforced}
(Section~\ref{sec:governance}). Ungoverned arms submitted 166 mutating
statements across 26 tasks and corrupted the warehouse twice; governed arms
submitted none and never attempted one, so the validator behind the rules
never fired.

The non-finding: on the 38, 30, 38 and 38 hard tasks outside the domain the
contract describes, it holds no detectable advantage over the prompt, and on
the strongest model the bare schema leads that slice. The gain is
domain-specific by construction. That is what a domain contract should do,
and it is the sharpest limit on what one benchmark can show.

\textbf{What a practitioner should take from this.} Spend on the
semantics first: write down what the columns mean and how the metrics are
defined, once per domain, in a form an agent can fetch one rule at a time.
That step bought most of the gain on every model, and the prose often
exists already. Build the retrieval scaffolding second; on frontier models
it is worth 14 to 15 points, on flash models close to nothing. Pre-compute
a macro only for a metric the agent demonstrably fails to derive, because
the shortfall macros close is already 5 points on the strongest model and
falling with capability, while the per-metric cost of writing and
maintaining them falls with nothing.
